% !TeX root = paper.tex
%
% PREAMBLE.TEX - LaTeX Package Configuration and Formatting Settings
%
% This file contains all package imports, formatting configurations, and custom
% command definitions for the academic paper template. It is loaded by paper.tex.
%
% IMPORTANT: Package loading order matters! Some packages must be loaded before
% others to avoid conflicts. The order below has been carefully arranged.
%
% Sections in this file:
%   - Document Encoding
%   - Bibliography Setup (using natbib)
%   - Page Format and Spacing
%   - Additional Packages
%   - Todo Notes System
%   - Captions and Figure/Table Settings
%   - TikZ Graphics
%   - Fonts and Typography (including microtype)
%   - Colors and Theorem Environments
%   - Section Formatting (titlesec - must load before hyperref)
%   - Hyperref and Cross-Referencing (cleveref loads after hyperref)
%   - Title Page Formatting
%   - Custom Backreference System
%   - User-Defined Commands

%% DOCUMENT ENCODING ===================================================
\usepackage[utf8]{inputenc}
\usepackage[T1]{fontenc}


%% BIBLIOGRAPHY ========================================================
\usepackage{csquotes}
\usepackage{xpatch}
\usepackage{natbib}

%% PAGE FORMAT =========================================================
\usepackage{sectsty} 			
\sectionfont{\Large}
\subsectionfont{\large}

\usepackage[verbose]{geometry}
\geometry{letterpaper}
\geometry{margin=2.5cm}

\setlength{\parskip}{0.5\baselineskip}	% change line skip after a paragraph finishes

%\usepackage{parskip}

\usepackage{setspace}
\onehalfspacing
%\doublespacing

% FOOTNOTES FORMAT =====================================================
%\usepackage[bottom,hang]{footmisc}
\usepackage[marginal]{footmisc}
\setlength{\footnotesep}{0.65\baselineskip}
\setlength{\footnotemargin}{-1.5ex}


%% ADDITIONAL PACKAGES =================================================
\usepackage{graphicx}
\graphicspath{{./figures/}} 

\usepackage{array} 
\usepackage{paralist} 
\usepackage{verbatim} 
\usepackage{amsmath}
\usepackage{amsthm}
\usepackage{booktabs}
\usepackage{multirow}
\usepackage{bigstrut}
\usepackage{amssymb}
\usepackage{bbm}
\usepackage{soul}
\usepackage{pdflscape}
\usepackage{multicol}
\usepackage{enumitem}
\usepackage[linewidth=1pt]{mdframed}
\usepackage{xspace}
\usepackage{tocloft}
\usepackage{etoc}
\usepackage{adjustbox}
\usepackage{anyfontsize}
\usepackage{ifdraft}
\usepackage{eucal}
\usepackage{rotating}
\usepackage{fancyvrb}
\usepackage{comment}
\usepackage{float}

% TODONOTES PACKAGE ====================================================
\usepackage[obeyFinal, colorinlistoftodos, textsize=small]{todonotes}
\setuptodonotes{fancyline, color=yellow, inline}

% https://mirror.hmc.edu/ctan/macros/latex/contrib/todonotes/todonotes.pdf
\newcounter{todoListItems}
\newcommand{\todoLink}[2][ ]{
% Increment counter
\addtocounter{todoListItems}{1}
\todo[%
caption={\protect\hypertarget{todo\thetodoListItems}{}Translation},
#1]
{
#2 \hfill
\hyperlink{todo\thetodoListItems}{$\uparrow$}
}
}

\newcommand{\td}{\todoLink}

\newcommand{\forAll}[2]{\td[color=yellow,inline,caption={\thetodoListItems:
    @All -- #1}]{\textbf{@All} -- #2}}
\newcommand{\forA}[2]{\td[color=green!40,inline,caption={\thetodoListItems:
    @A -- #1}]{\textbf{@A} -- #2}}
\newcommand{\forB}[2]{\td[color=blue!40,inline,caption={\thetodoListItems:
    @B -- #1}]{\textbf{@B} -- #2}}


% CAPTIONS =============================================================
\usepackage[labelfont=bf,justification=centering, position=top]{caption}
\usepackage[labelfont=bf, justification=centering, labelformat=simple]{subcaption}
\renewcommand\thesubfigure{(\alph{subfigure})}


\usepackage{pstricks}
%\usepackage[multiple]{footmisc}
\usepackage{etoolbox,xstring,xspace}

% FIXED COLUMN WIDTHS AND A WAY TO HIDE COLUMNS ========================
\newcommand{\PreserveBackslash}[1]{\let\temp=\\#1\let\\=\temp}
\newcolumntype{C}[1]{>{\PreserveBackslash\centering}p{#1}}
\newcolumntype{R}[1]{>{\PreserveBackslash\raggedleft}p{#1}}
\newcolumntype{L}[1]{>{\PreserveBackslash\raggedright}p{#1}}
\newcolumntype{H}{>{\setbox0=\hbox\bgroup}c<{\egroup}@{}}

% TIKZ PROGRAMMING =====================================================
\usepackage{tikz}
\usetikzlibrary{matrix,shapes,arrows,fit,tikzmark}
\usetikzlibrary{shapes.misc}
\usetikzlibrary{backgrounds}
\usetikzlibrary{shapes.multipart}
\usetikzlibrary{arrows.meta}
\usetikzlibrary{shapes.geometric}

%% CHANGE FONTS ========================================================
\usepackage[looser]{newpxtext} % LOOSER = BETTER WORD SPACE
\usepackage{textcomp}
\usepackage[upint]{newpxmath}
\usepackage{bm}

%% MICROTYPOGRAPHY =====================================================
% Improves typography and reduces overfull hbox warnings
\usepackage{microtype}

%% COLORS ==============================================================
\definecolor{ash}{RGB}{50,50,50}
\definecolor{darkred}{rgb}{0.35, 0.0, 0.0}
\definecolor{darkblue}{rgb}{0.0, 0.0, 0.5}
\definecolor{navyblue}{rgb}{0.0, 0.0, 0.8}
\definecolor{cadmiumgreen}{rgb}{0.0, 0.6, 0.24}
\definecolor{red_tol}{HTML}{D92120}
\definecolor{blue_tol}{HTML}{404096}
\definecolor{blugreen_tol}{HTML}{529DB7}
\setulcolor{darkred}

\newtheorem{theorem}{Theorem}
\newtheorem{proposition}{Proposition}
\newtheorem{lemma}{Lemma}
\newtheorem{definition}{Definition}

\newtheorem{requirement}{Requirement}
\newtheorem{assumption}{Assumption}
\newtheorem{corollary}{Corollary}

\theoremstyle{remark}
\newtheorem{example}{Example}

%\include{listings-stata}
\usepackage{xcolor}
\usepackage{tcolorbox}

% SECTION FORMATTING (must load before hyperref) ==================
\usepackage{titlesec}
%\titleformat{\paragraph}{\normalfont\normalsize\bfseries}{\theparagraph}{0pt}{}
\titlespacing*{\paragraph}{0pt}{1ex}{1em}
%\let\oldparagraph=\paragraph
%\renewcommand\paragraph[1]{\oldparagraph{#1.}}

%\titlespacing\section{0pt}{0.5ex}{0.5ex}
%\titlespacing\subsection{0pt}{0.5ex}{0.5ex}
%\titlespacing\subsubsection{0pt}{0.5ex}{0.5ex}


%% HYPERREF (should load near the end) =================================
\usepackage{hyperref}
\hypersetup{
    pdfstartview={FitH}, % fits the width of the page to the window
    pdftitle={XXX}, % title
    pdfauthor={Diego Jimenez Hernandez}, % author
    colorlinks=true,  % false: boxed links; true: colored links
    linkcolor=darkblue, % color of internal links (change box color with linkbordercolor)
    citecolor=darkred, % color of links to bibliography
    filecolor=magenta, % color of file links
    urlcolor=darkred, % color of external links
}

% PACKAGE FOR TITLE PAGE =====================================
\usepackage{titling}
\setlength{\droptitle}{-4em}              % PUSH TITLE UPWARDS
\pretitle{\begin{center}\huge\textbf}     % MAKE TITLE HIGHE AND BOLD
\posttitle{\par\end{center}\vskip 0.5em}  % FINISH THE TITLE AND ADD SPACE
%\renewcommand{\thanksscript}[1]{}        % REMOVE THANKS MARK FROM FOOTNOTE
%\renewcommand{\makethanksmarkhook}{\fontsize{9}{11}\selectfont}

% BACKREF FOR FIGURES AND TABLES =============================
\makeatletter

\AtBeginDocument{%
\let\origref\ref
\renewcommand*\ref[1]{%
  \origref{#1}\xlabel{#1}}
}
\newrobustcmd*\xlabel[1]{%
   \ifcsdef{siteref@doc@#1}{}{\csgdef{siteref@doc@#1}{,}}%
    \@bsphack%
    \begingroup
       \csxdef{siteref@doc@#1}{\csuse{siteref@doc@#1},\thepage}%
         \protected@write\@auxout{}%
        {\string\SiteRef{siteref@#1}{\csuse{siteref@doc@#1}}}%
     \endgroup
     \@esphack%
}

\newrobustcmd*\SiteRef[2]{\csgdef{#1}{#2}}

\newrobustcmd*\xref[1]{%
\ifcsundef{siteref@#1}{%
     \@latex@warning@no@line{Label `#1' not defined}
     }{%
    \begingroup
      \StrGobbleLeft{\csuse{siteref@#1}}{2}[\@tempa]\relax%
      \def\@tempb{}%
      \@tempcnta=0\relax%
      \@tempcntb=\@ne\relax%
      \def\do##1{\advance\@tempcnta\@ne}%
      \expandafter\docsvlist\expandafter{\@tempa}%
       \def\do##1{%
         \ifnum\@tempcntb=\@tempcnta\relax%
            \hyperpage{##1}%
         \else
            \hyperpage{##1},%
          \fi%
          \advance\@tempcntb\@ne
       }%
       [\expandafter\docsvlist\expandafter{\@tempa}]\xspace%
    \endgroup
   }%
}
\makeatother

% USER-WRITTEN COMMAND DEFINITIONS ===========================
% BACKREFERENCE COMMAND
\newcommand{\backreference}[1]{%
\ifoptionfinal
    {}%
    { 

    \noindent
    {\footnotesize\textit{Backreferenced:} 
     \texttt{\detokenize{#1}} \xref{#1}}}

    \noindent%
}

% OUTLINE
\newcommand{\outline}[1]{%
\begin{tcolorbox}%
\footnotesize%
	#1
\end{tcolorbox}
}

% RED TEXT
\newcommand{\redtext}[1]{\textcolor{red}{#1}}

% XXX TO FILL OUT
\newcommand{\xx}{\textcolor{magenta}{\texttt{<XX>}}}

% TAKEAWAY COMMAND
\newcommand{\takeaway}[1]{\begin{center}\textit{#1}\end{center}}


% SINGLE SPACING IN TABLES AND FIGURES
\AtBeginEnvironment{table}{\vspace{-0.5\baselineskip}\singlespacing}
\AtBeginEnvironment{figure}{\vspace{-0.5\baselineskip}\singlespacing}

\AtEndEnvironment{table}{\vspace{-0.5\baselineskip}}
\AtEndEnvironment{figure}{\vspace{-0.5\baselineskip}}








